% Introduction
%
% Suggested flow for a CS review/survey paper:
% 1. Broad context: why software testing matters
% 2. Narrowing: what is changing ("next generation")
% 3. Problem statement / gap this review addresses
% 4. Research questions (RQs)
% 5. Contributions (bulleted)
% 6. Paper structure (roadmap)

\IEEEPARstart{S}{oftware} testing remains one of the most resource-intensive
activities in the software development lifecycle, commonly consuming
30--50\% of overall project effort~\cite{placeholder2024}. Traditional
testing techniques -- manual test case design, coverage-driven unit
testing, and scripted regression suites -- are increasingly challenged
by the scale, complexity, and release velocity of modern software
systems, including microservices, machine-learning-enabled components,
and continuously deployed cloud-native applications.

In response, a new generation of testing approaches has emerged,
leveraging artificial intelligence (AI), machine learning (ML),
large language models (LLMs), self-adaptive test oracles, and
autonomous exploratory agents. These approaches promise to reduce
manual effort, improve defect detection, and adapt testing strategies
dynamically. However, the research landscape is fragmented: studies
differ in terminology (e.g., ``intelligent testing,'' ``autonomous
testing,'' ``AI-augmented testing''), evaluation benchmarks, and
maturity levels, making it difficult for researchers and practitioners
to assess the state of the art.

\subsection{Problem Statement}
No comprehensive, up-to-date systematic review consolidates these
next-generation software testing techniques under a unified
classification framework with rigorous methodological screening.

\subsection{Research Questions}
This review addresses the following research questions (RQs):
\begin{itemize}
  \item \textbf{RQ1:} What next-generation software testing techniques
        have been proposed in the literature, and how can they be
        classified?
  \item \textbf{RQ2:} What levels of automation and autonomy do these
        techniques achieve across the test life cycle (design,
        generation, execution, oracle, maintenance)?
  \item \textbf{RQ3:} What empirical evidence exists regarding the
        effectiveness, efficiency, and industrial adoption of these
        techniques?
  \item \textbf{RQ4:} What are the open challenges and research gaps
        for next-generation software testing?
\end{itemize}

\subsection{Contributions}
\begin{itemize}
  \item A systematic, reproducible literature review protocol
        (Section~\ref{sec:methodology}) covering \textit{[DATE\_RANGE]}.
  \item A taxonomy of next-generation software testing techniques
        derived from \textit{[N\_FINAL]} primary studies.
  \item A synthesis of empirical evidence on effectiveness and
        adoption barriers.
  \item A forward-looking research agenda (Section~\ref{sec:conclusion}).
\end{itemize}

\subsection{Paper Structure}
The remainder of this paper is organized as follows.
Section~\ref{sec:background} provides background and motivation.
Section~\ref{sec:methodology} details the review protocol.
Section~\ref{sec:related-work} discusses related reviews and surveys.
Section~\ref{sec:results} presents the synthesis of findings.
Section~\ref{sec:discussion} discusses implications.
Section~\ref{sec:validity} addresses threats to validity, and
Section~\ref{sec:conclusion} concludes with future directions.
